% ------------------------------------------------------------------------
% file `3_uaa1_fonctions-ineqations-systeme_02_2023-2024-transferer-10-exercise-body.tex'
%   in folder `xsim/'
%
%     exercise of type `transferer' with id `10'
%
% generated by the `transferer' environment of the
%   `xsim' package v0.21 (2022/02/12)
% from source `3_uaa1_fonctions-ineqations-systeme_02_2023-2024' on 2023/10/20 on line 261
% ------------------------------------------------------------------------
    \tcbsidebyside[blanker, sidebyside adapt=left]{
        \includegraphics[width=5cm]{tikz/trapeze.tikz}
    }{
        $TRAP$ est un trapèze rectangle en $A$ et en $P$ tel que :\\
        $|TP| = \SI{3}{cm}$, $|PA| = \SI{5}{cm}$ et $|AR|=\SI{4}{cm}$.\medbreak
        $M$ est un point mobile du segment $[PA]$ et on note $x$ la longueur du segment $[PM]$ en cm.
    }
    \begin{tasks}
        \task Montre, grâce au dessin, que l'aire du triangle $PTM$ est $1,5x$ et que l'aire du triangle $ARM$ est $-2x+10$.
    \end{tasks}

    Les droites ci-dessous sont les représentations graphique des fonctions $f$ et $g$ qui associent respectivement \(x\) à l'aire des triangles.

    \includegraphics[width=\linewidth]{tikz/graphe_trapeze.tikz}

    \begin{tasks}[resume=true]
        \task Pour quelles valeur de $x$ l'aire du triangle $ARM$ est-elle égale à \qty{8}{cm^2} ?
        \task Lorsque $x$ est égal à \qty{4}{cm}, quelle est l'aire du triangle $ARM$ ?
        \task Montre par calcul que la valeur de $x$ pour laquelle les deux aires sont égales est $\dfrac{20}{7}$.
    \end{tasks}
